% =====================================================================
% 別冊 付録A おすすめのブランチ運用ルールとPRの書き方
% 本文（第1章）で軽く触れた内容の、実践的なテンプレート集。
% =====================================================================
\chapter{おすすめのブランチ運用ルールとPRの書き方}

第1章では、ブランチの切り方やPRの作り方を扱いました。この付録では、\textbf{チームで気持ちよく開発するための「型」}を、そのまま真似できる形でまとめます。絶対の決まりではありませんが、\textbf{迷ったらこれに沿えばまず問題ありません}。チームで最初に決めておくと、後の混乱を防げます。

% ---------------------------------------------------------------------
\section{おすすめのブランチ運用ルール}

\begin{center}
\begin{tabularx}{\linewidth}{@{}>{\raggedright\arraybackslash}p{6.6cm} X@{}}
\toprule
ルール & なぜそうするか \\
\midrule
\cmd{main} は直接編集しない & \cmd{main} を常に「動く・完成した状態」に保つため。GitHubの設定で \cmd{main} への直接 \cmd{push} を禁止することもできる \\
\addlinespace[3pt]
1タスク＝1ブランチ & 目的が1つに絞られ、レビューとマージがきれいにまとまる \\
\addlinespace[3pt]
作業前に最新の \cmd{main} から枝を出す & 古い \cmd{main} から分岐すると、後でコンフリクトが起きやすい \\
\addlinespace[3pt]
こまめにコミット・こまめに \cmd{push} & 事故ってもすぐ戻せる。作業内容を他のメンバーが確認できる \\
\addlinespace[3pt]
PRは小さく保つ & 変更が小さいほどレビューが速く、正確になる \\
\addlinespace[3pt]
マージが済んだらブランチを削除する & 使わないブランチが増えて散らかるのを防ぐ \\
\bottomrule
\end{tabularx}
\end{center}

\medskip

\subsection{ブランチ命名のおすすめ}

ブランチ名は \textbf{\cmd{種類/内容}} の形（英小文字・単語はハイフン区切り）にすると、一覧で目的がひと目で分かります。

\begin{center}
\begin{tabularx}{\linewidth}{@{}>{\raggedright\arraybackslash}p{2.6cm} >{\raggedright\arraybackslash}p{5.4cm} X@{}}
\toprule
接頭辞 & 用途 & 例 \\
\midrule
\cmd{feature/} & 新機能の追加 & \cmd{feature/add-login} \\
\cmd{fix/} & バグ修正 & \cmd{fix/typo-in-readme} \\
\cmd{docs/} & ドキュメントの変更 & \cmd{docs/update-setup} \\
\cmd{refactor/} & 動作を変えない整理 & \cmd{refactor/tidy-utils} \\
\cmd{chore/} & 設定・雑務 & \cmd{chore/update-gitignore} \\
\bottomrule
\end{tabularx}
\end{center}

\medskip

% ---------------------------------------------------------------------
\section{伝わるPRの書き方}

PRは、レビュアーに\textbf{「何を・なぜ変えたか」}を伝える場です。差分（コード）を見れば「どう変えたか」は分かりますが、\textbf{「なぜ」}は書かないと伝わりません。次の点を意識すると、レビューが速く進みます。

\begin{itemize}[leftmargin=1.4em]
  \item \textbf{タイトルは一目で分かるように}。おすすめは \cmd{feat: ログイン機能を追加} のように\textbf{種類を頭に付ける形}（\cmd{feat} 新機能／\cmd{fix} 修正／\cmd{docs} 文書 …）
  \item \textbf{説明には「なぜ」と「背景」を書く}。差分だけでは分からない意図を補う
  \item \textbf{1つのPRは1つの目的に絞る}。ついでの修正を混ぜない
  \item \textbf{動作確認したことを書く}。何をどう確かめたかがあると、レビュアーは安心できる
  \item \textbf{まだ作業中なら Draft（下書き）PR にする}。完成したら「Ready for review」に切り替える
\end{itemize}

% ---------------------------------------------------------------------
\section{PRテンプレート}

毎回ゼロから説明を書くのは大変です。\textbf{PRテンプレート}を用意しておくと、PR作成時に説明欄へ\textbf{自動で見出しが挿入}され、埋めるだけで済みます。

リポジトリの一番上の階層に \cmd{.github/PULL\_REQUEST\_TEMPLATE.md} というファイルを置くと、GitHubがPR作成時に自動で読み込みます。次はそのまま使えるテンプレート例です。

\begin{codebox}[text]
\begin{code}
## 概要
<!-- このPRで何をしたかを1〜2行で -->

## 変更内容
-

## 変更の理由・背景
<!-- なぜこの変更が必要か。差分だけでは分からない意図を書く -->

## 動作確認
- [ ] ビルド / テストが通ることを確認した
- <!-- 確認した手順・結果を書く -->

## 関連 Issue
<!-- あれば #番号。なければ「なし」 -->

## レビューで見てほしい点
<!-- 特に確認してほしい箇所があれば -->
\end{code}
\end{codebox}

テンプレート内の記法について補足します。\cmd{- [ ]} はGitHub上で\textbf{チェックボックス}として表示され、クリックでチェックできます。\cmd{<!-- ... -->} はコメントで、\textbf{記入後は消して}構いません（残しても表示はされません）。

% ---------------------------------------------------------------------
\section{記入例}

上のテンプレートを、\textbf{ログイン機能を追加するPR}で埋めた例です。

\noindent\textbf{タイトル}

\begin{codebox}[text]
\begin{code}
feat: ログイン機能を追加
\end{code}
\end{codebox}

\noindent\textbf{説明}

\begin{codebox}[text]
\begin{code}
## 概要
メールアドレスとパスワードでログインできる画面を追加した。

## 変更内容
- ログイン画面（login.html）を追加
- 入力チェックと認証処理を実装
- ログイン後にトップページへ遷移する処理を追加

## 変更の理由・背景
会員限定ページを用意するため、まず入口となるログイン機能が必要だった。

## 動作確認
- [x] ビルド / テストが通ることを確認した
- 正しいメール・パスワードでログインできることを確認
- 誤ったパスワードでエラーメッセージが出ることを確認

## 関連 Issue
#12

## レビューで見てほしい点
パスワードの扱い方（平文で残していないか）を重点的に見てほしい。
\end{code}
\end{codebox}

\begin{hintbox}
最初から完璧を目指す必要はありません。\textbf{「概要」と「変更の理由」だけでも書いてあると、レビュアーの理解がぐっと速くなります}。チームで運用しながら、自分たちに合う項目に育てていきましょう。
\end{hintbox}
